\section{Economic action contrasts and independent risk assessment}\label{sec:r17-contrast-proofs}
This section concerns the finite continuation economy defined in
Section~\ref{sec:r16menumain}. All expectations below use its unclipped Gaussian
law and its stored economic coefficients. Conditioning on a fitted continuation
is different from assuming a fixed feature dictionary before training. The
results in this section permit nonlinear fitting, data-dependent features, and
actions chosen by the fitted critic, provided the assessment innovations are
independent of these choices.

Let $\mathcal F$ contain the complete training record, a finite collection of
queries indexed by $j$, two feasible actions at each query, and every deployed
prediction to be assessed. Let $x_j^+$ and $x_j^0$ denote their postdecision
states. Write
\[
 c_j=S(x_j^+)-S(x_j^0),\qquad
 p_{Nj}=\widehat S_N(x_j^+)-\widehat S_N(x_j^0),\qquad p_{Rj}\in\mathbb R.
\]
Here $S$ is the true continuation of the reference policy in the finite economy;
$p_{Rj}$ may be a realized cached simulation predictor. Conditional on
$\mathcal F$, all these quantities are fixed. The risk uses the same action pair
and the same query weights for both predictors:
\begin{equation}\label{eq:r17-risk}
 \mathcal R_M=\frac1J\sum_{j=1}^J(p_{Mj}-c_j)^2,
 \qquad M\in\{N,R\}.
\end{equation}
This definition does not compare the payoffs of two different actors. It does
not average over unrealized Raw caches or unexecuted training seeds.

\begin{proposition}[Action contrasts and economic decisions]\label{prop:r17-decisions}
Let $r_j$ be the known current-reward difference between the two actions.
The selector that chooses the first action when $r_j+p_{Mj}\geq0$ has loss,
relative to the better of these two actions, bounded by $|p_{Mj}-c_j|$.
Consequently its average loss is at most $\sqrt{\mathcal R_M}$. An additive
error of at most $\eta_j$ in maximizing the predicted two-action objective adds
$\eta_j$ to the pointwise bound. For a finite action catalog with reference
contrast $c_{j0}=p_{Mj0}=0$, an $\eta_j$-approximate predicted maximizer instead
has loss at most
\[
 2\max_a|p_{Mja}-c_{ja}|+\eta_j.
\]
If the catalog risk at query $j$ is $\sum_a w_{ja}(p_{Mja}-c_{ja})^2$ and
$w_{ja}\geq\underline w>0$ for every action, the last bound is at most twice
the square root of this risk divided by $\underline w$, plus $\eta_j$.
\end{proposition}
\begin{proof}
Put $v=r_j+c_j$ and $\widehat v=r_j+p_{Mj}$. When the selectors based on $v$
and $\widehat v$ agree, the loss is zero. When they disagree, zero lies between
$v$ and $\widehat v$, so the loss $|v|$ is at most $|v-\widehat v|$.
For an approximate choice, the same argument bounds a wrong choice by
$|v-\widehat v|+\eta_j$. Averaging and applying Cauchy--Schwarz proves the
risk bound. For a catalog let $a^*$ maximize the true objective and let
$\widehat a$ be the approximate predicted maximizer. Insert the predicted
objective at $a^*$ and $\widehat a$ into their true difference; the middle
term is at most $\eta_j$ and each of the other terms is at most the maximum
contrast error. Finally, each squared error is at most the weighted risk
divided by its weight. This proves every assertion.
\end{proof}

The first assertion provides an economic unit for the risk without asserting
that risk ordering determines payoff ordering. Such an assertion is false:
when the true net contrast is $.01$, predictions $-.01$ and $.2$ have squared
errors $.0004$ and $.0361$, respectively, but only the less accurate predictor
chooses the better action. A risk comparison and an actual payoff comparison
therefore remain separate empirical results. The two-action bound is not a
certificate over all vector actions or over the full adapted diffusion class.

\begin{proposition}[Cancellation of common continuation noise]\label{prop:r17-cancellation}
Conditional on $\mathcal F$, let $Z_{jk}$ be independent across queries and
replicates, with finite second moments and $\mathbb E[Z_{jk}\mid\mathcal F]=c_j$.
Put $d_j=p_{Nj}-p_{Rj}$ and
\[
 T_{jk}=d_j(p_{Nj}+p_{Rj}-2Z_{jk}).
\]
Then the uniform average of $T_{jk}$ is unbiased for
$\mathcal R_N-\mathcal R_R$, and
\begin{equation}\label{eq:r17-variance}
 \operatorname{Var}(T_{jk}\mid\mathcal F)
   =4d_j^2\operatorname{Var}(Z_{jk}\mid\mathcal F).
\end{equation}
An identical pair of predictions has identically zero comparison statistic.
Neither an independent estimate of $c_j^2$ nor two independent assessment
banks are required for this risk difference. For absolute risk, two independent
replicates $Z_{j1},Z_{j2}$ give the unbiased statistic
$(p_{Mj}-Z_{j1})(p_{Mj}-Z_{j2})$.
\end{proposition}
\begin{proof}
The algebraic identity
\[
 (p_{Nj}-z)^2-(p_{Rj}-z)^2
   =(p_{Nj}-p_{Rj})(p_{Nj}+p_{Rj}-2z)
\]
removes the common $z^2$ term. Take conditional expectations and average.
The only random component of the right side is $-2d_jZ_{jk}$, proving
\eqref{eq:r17-variance}. For the last assertion, conditional independence
factors the expectation of the product into
$(p_{Mj}-c_j)^2$. Reusing one replicate instead would add its conditional
variance. No assumption about how the predictions were fitted enters this
argument, because fitting is included in $\mathcal F$.
\end{proof}

\begin{corollary}[An economic contrast training objective]\label{cor:r17-training}
For fixed training-query weights, the conditional expectation of
$\sum_j w_j\{f(x_j^+)-f(x_j^0)-Z_j\}^2$ equals the corresponding true contrast
risk plus $\sum_jw_j\operatorname{Var}(Z_j\mid\mathcal F)$. The latter does
not depend on $f$. Thus minimizing the population loss over scalar functions
minimizes contrast risk, while preserving the common scalar continuation and
its path-consistent differences. The loss is invariant to an additive constant
in $f$. A fresh assessment bank is required after empirical fitting or model
selection; a training loss is not an independent risk certificate.
\end{corollary}
\begin{proof}
Expand the square around the conditional mean $c_j$. The cross term has zero
conditional mean. Summing gives the identity. Additive constants cancel in
each difference, and differences of a single scalar function telescope around
every finite cycle. These are algebraic properties, not statistical claims.
\end{proof}

\begin{proposition}[Centered finite-sample comparison]\label{prop:r17-risk-certificate}
Fix $q_j,b_j,t_j,e_j$ before assessment, with $b_j,t_j,e_j\geq0$. Suppose the
exact residual $U_{jk}=Z_{jk}-q_j$ obeys
$\mathbb E(|U_{jk}|-b_j)_+\leq t_j$. Suppose the executed residual
$\widetilde U_{jk}$ can be coupled with $U_{jk}$ so that
$\mathbb E|\widetilde U_{jk}-U_{jk}|\leq e_j$.
Define the known offset and the bounded random statistic by
\[
 K_j=d_j(p_{Nj}+p_{Rj}-2q_j),\qquad
 W_{jk}=-2d_j\operatorname{clip}_{b_j}(\widetilde U_{jk}).
\]
Then $|W_{jk}|\leq B:=\max_j2|d_j|b_j$ and the difference between the
risk contrast and $J^{-1}\sum_jK_j+\mathbb E\overline W$ is at most
$\varepsilon=J^{-1}\sum_j2|d_j|(e_j+t_j)$.
For $n=JK\geq2$ independent observations with $K$ replicates per query, a
simultaneous two-sided interval for one comparison is
\begin{equation}\label{eq:r17-eb}
 \frac1J\sum_jK_j+\overline W
 \ \mathbin{\pm}\left\{
 \sqrt{\frac{2s_W^2\log(4/\alpha_*)}{n}}
 +\frac{14B\log(4/\alpha_*)}{3(n-1)}+\varepsilon\right\},
\end{equation}
with coverage at least $1-\alpha_*$. For $H$ prespecified comparisons,
$\alpha_*\leq\alpha/H$ yields simultaneous coverage at least $1-\alpha$.
The observations need not be identically distributed.
\end{proposition}
\begin{proof}
Clipping is one-Lipschitz. The triangle inequality gives
\[
 \mathbb E|U_{jk}-\operatorname{clip}_{b_j}(\widetilde U_{jk})|
 \leq t_j+e_j.
\]
Multiply by $2|d_j|$ and average to obtain the expectation allowance. The range
is immediate. For $B>0$, transform $W$ from $[-B,B]$ to $[0,1]$, apply
\citet[Theorem~11]{maurer2009} to each tail with probability
$\alpha_*/2$, and transform back. Its sample variance is the usual unbiased
sample variance; the theorem allows independent nonidentical observations.
The union bound proves the family assertion. For $B=0$ the statistic is
identically zero and only the known offset and deterministic allowance remain.
\end{proof}

The implementation takes $q_j$ to be a representative of the known terminal
quadratic contrast. The inherited finite-economy geometric proof supplies
$b_j,t_j$ before observation; its protected-arithmetic and Gaussian-clipping
accounts supply $e_j$. Antithetic signs are averaged before they enter the
sample: they are not counted as independent draws. Outward interval arithmetic
also encloses the known offset, coefficient subtraction, coefficient
multiplication, and final endpoint addition. The sample moments are calculated
as exact rational moments of the stored binary64 observations before directed
Decimal evaluation of the radius. Identical predictors are assigned their
exact zero identity, not an artificially positive comparison width.
